\documentclass{zhvt-classic}

% 使用本地源流明体（GenRyuMin2 TW）按文件路径加载，不依赖系统字体名
% 版框（内框／界栏／书口侧栏／鱼尾）改用与版心印章同一印泥棕色：
%   印章图片 yshs.png 实色 = RGB(135,68,52) = HTML 874434
\usepackage{xcolor}
\usepackage{tikz}
% sealbrown 双注册：xcolor 供 TikZ/pgf，L3 供 cls 版框
\definecolor{sealbrown}{HTML}{874434}
\ExplSyntaxOn
\color_set:nnn { sealbrown } { HTML } { 874434 }
\ExplSyntaxOff

\zhvtset{
  grid color       = sealbrown,
  main font        = GenRyuMin2TW-SB,
  main font path   = C:/Users/Administrator/texmf/fonts/opentype/public/GenRyuMin2TW/,
jiazhu font      = GenRyuMin2TW-M,
  jiazhu font path = C:/Users/Administrator/texmf/fonts/opentype/public/GenRyuMin2TW/,
  % 页码框容纳 4 字：正常页「九十七」+前置空 1 格=4 行恰好；
  % 100+ 页「一百零一」共 5 行超出，末字「一/二/三」被框裁掉（2026-10-04 实测），故放宽到 5
  shukou page number chars = 5,
  % 书口书名框高从默认 6 字减到 4 字，把省出的 2 字高拨给 chapter 框，
  % 使 6 字书名+卷名放进 9 字高的 chapter 框后不再压线（2026-10-05）
  shukou title chars = 4,
}

\title{史记}[商君列传]
\author{汉·司马迁}
\maker{闲叟漫记·重製}


% ===== 索隐述赞样式（2026-11-08 四次定稿 · 愚大人亲手改定参数）=====
% ⚠ 竖排前提：自定义字体家族必须带 RawFeature={vertical:+vert}；
%   simhei.ttf 加 +vhal 会整体加载失败静默回退，只写 +vert；simfang 同理。
\newCJKfontfamily{\shuzanfont}[Path=C:/Windows/Fonts/,RawFeature={vertical:+vert}]{simhei.ttf}
\newCJKfontfamily{\shuzanbody}[Path=C:/Windows/Fonts/,RawFeature={vertical:+vert}]{simfang.ttf}
% 题头：152x30pt 圆角 6pt 棕底白字；baseline (0,2) 使整框落入界行格子不压版框线；
% 白字节点 x 不补偿、y +8.2pt（像素级墨迹实测 dx/dy≈0，L/R 留白约 20pt）。
\newcommand{\shuzanhead}[1]{%
  \begin{tikzpicture}[baseline={(0,2)}]
    \fill[sealbrown, rounded corners=6pt]
      (-76pt,-15pt) rectangle (76pt,15pt);
    \node[text=white, inner sep=0pt, anchor=center] at (0pt,8.2pt)
      {\shuzanfont\fontsize{23}{23}\selectfont #1};
  \end{tikzpicture}%
}
% ★ 修正书口卷名位置：chapter 框起始位置改回 cls 默认值，避免全局影响。
%   只在每个卷首页通过 \juan 命令里的空命令微调该页位置（安全，不会跨页污染）。
\ExplSyntaxOn
\AtBeginDocument{
  \dim_set:Nn \l__zhvt_shukou_chapter_title_start_dim
    { \l__zhvt_shukou_title_height_dim + \l__zhvt_font_size_dim * 2 }
  \dim_set:Nn \l__zhvt_shukou_page_number_start_dim
    { \l__zhvt_shukou_chapter_title_start_dim + \l__zhvt_shukou_chapter_title_height_dim }
}
\ExplSyntaxOff

\begin{document}

% ★★ 版心（书口）排布覆写：把书名放到鱼尾之下（2026-10-04）
%
%   cls 默认版心自上而下：书名（title 静态框，在鱼尾之上）→ 鱼尾 → 篇名（chapter 动态框）→ 页码
%   本模板改为：          鱼尾（位置不动）→ 书名 → 篇名 → 页码分割线 → 页码
%
%   手法：把 cls 放在 title 静态框里的书名，改放进 chapter 动态框（鱼尾之下）；
%         title 静态框留空。书名取 \l__zhvt_title_str，即 \title{...} 的第一参数，
%         所以换书时只改 \title，不必在此重复写书名。
%
%   ⚠ 三条必须遵守的写法（2026-10-04 实测）：
%     ① 汉字必须用 \hbox{...} 包住。flowfram 会把框内容装进一个宽 14.4mm（= 一个 line_sep
%        = 书口栏宽）的纵向 minipage 并按该宽度折行，而整页靠 PDF /Rotate 90 显示、
%        框内文字本身不变向 —— 不包 \hbox 时「围炉夜话」4 字会被逐字折成 4 行，
%        在版面上就成了横排 4 字并横向溢出版框。
%     ② 竖排下**一个 \hbox 就是一整列**。书名与篇名要同处一列，必须放进同一个 \hbox；
%        分成两个 \hbox 时，后一个会被当成新的一列、甩到隔壁栏去
%        （实测「序」跑到书口栏外侧 14.4mm 的正文列上）。中间的 \quad 在框内是横向空白，
%        显示出来正好是书名与篇名之间的纵向间隔。
%     ③ ★ 内容盒要写成 **与框等宽的定宽盒** `\hbox to \l__zhvt_shukou_chapter_title_height_dim{…\hfil}`。
%        实测：chapter 框是把内容盒沿纵向**居中**摆放的，内容宽度一变，书名就跟着漂：
%        自然宽度的 \hbox（书名 38.4mm）落在 x=157.2mm，加篇名后（57.6mm）落在 x=152.4mm，
%        目录页（无篇名）与正文页的书名就差半格。做成等宽盒后，无论有没有篇名，
%        内容固定从框顶起排，实测两页均落在 x=142.8mm，不再漂。
\ExplSyntaxOn
% 把书口 chapter 框（书名+篇名）起始位置下移 1 字高，避免 4 字卷名顶到鱼尾
% ⚠ cls 的 \__zhvt_calculate_dimensions: 在 \AtBeginDocument 里执行，会重算该 dim；
%   必须把覆写挂到 begin document 之后，否则会被二次覆盖（2026-10-05 实测）。
% (旧方案已移到导言区，直接修改 chapter 框起始位置)
% 卷标题命令 #1=书口篇名(鱼尾) #2=正文卷名（含卷序，卷名内空格用 ~）
%
% ★★ 目录页两项修订（2026-10-05，愚大人指出）★★
%   ① 页码起点统一对齐：
%      原状 —— 目录项写成 \par\hspace{4em}卷名\hspace{4em}页码，页码起点 = 卷名实宽 + 4em。
%      卷名实宽随字数变：「卷一~梁惠王上」6 字 +~、「卷七~离娄上」5 字 +~，
%      差 1 个字位 → 卷七~卷十 4 条的页码整体比其余 10 条高 27.21pt（一个字），看着上下不齐。
%      改法 —— 卷名补全角空格（宽 = 1 字 = 27.215pt）补足到 **7 字符位**再写进目录，
%      14 条页码起点遂全部落在同一条竖线上（起点对齐；末字随页码字数自然参差）。
%      ⚠ 只对目录里的副本 \l_tmpc_tl 补，正文卷名仍用 \l_tmpa_tl 原样 —— 正文像素不受影响。
%   ② 卷名整体上移半字位：前导 \hspace 由 4em 改 3.5em。
%      实测原状目录块在版框内偏下（天头留白 139.6pt / 地脚留白 110.7pt），
%      上移 0.5em = 13.6pt 后上下基本均分，视觉居中。
\NewDocumentCommand { \juan } { mm } {
	\clearpage
	\tl_set:Nn \l__zhvt_chapter_title_tl { #1 }
	\tl_set:Nn \l_tmpa_tl { #2 }
	\tl_set_eq:NN \l_tmpc_tl \l_tmpa_tl
	\int_step_inline:nn { 7 - \str_count:n { #2 } }
	  { \tl_put_right:Nn \l_tmpc_tl { 　 } }
	\tl_set:Nx \l_tmpb_tl { \zhnum { page } }
	\tl_put_left:Nn \l_tmpb_tl { \hspace { 4em } }
	\tl_put_left:Nn \l_tmpb_tl { \l_tmpc_tl }
	\tl_put_left:Nn \l_tmpb_tl { \par \hspace { 3.5em } }
	\addtocontents { toc } { \l_tmpb_tl }
	\skip_horizontal:n { \l__zhvt_font_size_dim * 4 }
	\l_tmpa_tl
	\par
	\thispagestyle { empty }
}
\cs_set_protected:Npn \__zhvt_set_shukou_content: {
	% ★ 鱼尾上方留白（title 静态框）嵌入闲叟钤印：
	%   整页 /Rotate 90，故此处把图用 \rotatebox{90} 转正为竖章；
	%   尺寸卡在版心条幅内（书内约 7mm 宽 × 33.8mm 高，留余量不压鱼尾）。
	\setstaticcontents * { title } {
		\vbox_to_ht:nn { \l__zhvt_baseline_skip_dim } {
			\vskip 14.2pt
			\hbox_to_wd:nn { \l__zhvt_shukou_title_height_dim } {
				\hfil
				\rotatebox [ origin = c ] { 90 } { \includegraphics [ width = 7mm , height = 33.8mm ] { yshs.png } }
				\hfil
			}
		}
	}
% ★ 版心（书口）统一改法（2026-11-16 愚大人定）：18pt＋框内居中＋书口文字同版框色
%   ① 原 27.2pt 时 7 字长篇名（如「廉颇蔺相如列传」）末字溢出 chapter 框、压住页码；
%      改 \fontsize{18pt}{18pt}，前后 \hfil 令字串在框内居中（不随篇名字数漂移）。
%   ② \vbox{\vskip 8.2pt} 定量下移，抵消缩字号后栏宽方向 8.2pt 上偏（像素实测 y 中心 543.3 = 原版）。
%      ⚠ \raisebox / \vspace* 在 flowfram 框内无效；换字号须同步换补偿：16pt→10.1、19pt→7.4。
%   ③ chapter 与 page~number 均置 \color{sealbrown}（HTML 874434），与版框／界行同色；正文仍黑。
%   ⚠ 页码字位、钤印、鱼尾、界行、正文一律未动；五部合订本同此口径。
	\setdynamiccontents * { chapter }
	{ \color{sealbrown} \vbox { \vskip 8.2pt \hbox to \l__zhvt_shukou_chapter_title_height_dim
		{ \hfil \fontsize{18pt}{18pt}\selectfont \l__zhvt_title_str \quad \l__zhvt_chapter_title_tl \hfil } } }
	\setdynamiccontents * { page~number } { \color{sealbrown} \quad \zhnum { page } }
}
\__zhvt_set_shukou_content:
\ExplSyntaxOff

\maketitle

% 单篇，无多卷目录，不渲染目录页

\juan{商君列传}{卷六十八~商君列传第八}
\hspace*{2em}商君者。卫之诸庶孽公子也。名鞅。姓公孙氏。其祖本姬姓也。鞅少好刑名之学。事魏相公叔痤为中庶子。公叔痤知其贤。未及进。会痤病。魏惠王亲往问病。曰。﹁公叔病有如不可讳。将柰社稷何。﹂公叔曰。﹁座之中庶子公孙鞅。年虽少。有奇才。原王举国而听之。﹂王嘿然。王且去。座屏人言曰。﹁王即不听用鞅。必杀之。无令出境。﹂王许诺而去。公叔痤召鞅谢曰。﹁今者王问可以为相者。我言若。王色不许我。我方先君後臣。因谓王即弗用鞅。当杀之。王许我。汝可疾去矣。且见禽。﹂鞅曰。﹁彼王不能用君之言任臣。又安能用君之言杀臣乎。﹂卒不去。惠王既去。而谓左右曰。﹁公叔病甚。悲乎。欲令寡人以国听公孙鞅也。既又劝寡人杀之。岂不悖哉。﹂\\
\hspace*{2em}公叔既死。公孙鞅闻秦孝公下令国中求贤者。将修缪公之业。东复侵地。乃遂西入秦。因孝公宠臣景监以求见孝公。孝公既见卫鞅。语事良久。孝公时时睡。弗听。罢而孝公怒景监曰。﹁子之客妄人耳。安足用邪。﹂景监以让卫鞅。卫鞅曰:﹁吾说公以帝道。其智不开悟矣。﹂后五日。复求见鞅。鞅复见孝公。益愈。然而未中旨。罢而孝公复让景监。景监亦让鞅。鞅曰。﹁吾说公以王道而未入也。请复见鞅。﹂鞅复见孝公。孝公善之而未用也。罢而去。孝公谓景监曰。﹁汝客善。可与语矣。﹂鞅曰。﹁吾说公以霸道。其意欲用之矣。诚复见我。我知之矣。﹂卫鞅复见孝公。公与语。不自知跶之前於席也。语数日不厌。景监曰。﹁子何以中吾君。吾君之驩甚也。﹂鞅曰。﹁吾说君以帝王之道比三代。而君曰。﹁久远。吾不能待。且贤君者。各及其身显名天下。安能邑邑待数十百年以成帝王乎。﹂故吾以彊国之术说君。君大说之耳。然亦难以比德於殷周矣。﹂\\
\hspace*{2em}孝公既用卫鞅。鞅欲变法。恐天下议己。卫鞅曰。﹁疑行无名。疑事无功。且夫有高人之行者。固见非于世。有独知之虑者。必见敖於民。愚者闇於成事。知者见於未萌。民不可与虑始而可与乐成。论至德者不和於俗。成大功者不谋於众。是以圣人苟可以彊国。不法其故。苟可以利民。不循其礼。﹂孝公曰。﹁善。﹂甘龙曰。﹁不然。圣人不易民而教。知者不变法而治。因民而教。不劳而成功。缘法而治者。吏习而民安之。﹂卫鞅曰。﹁龙之所言。世俗之言也。常人安於故俗。学者溺於所闻。以此两者居官守法可也。非所与论於法之外也。三代不同礼而王。五伯不同法而霸。智者作法。愚者制焉。贤者更礼。不肖者拘焉。﹂杜挚曰。﹁利不百。不变法。功不十。不易器。法古无过。循礼无邪。﹂卫鞅曰。﹁治世不一道。便国不法古。故汤武不循古而王。夏殷不易礼而亡。反古者不可非。而循礼者不足多。﹂孝公曰。﹁善。﹂以卫鞅为左庶长。卒定变法之令。\\
\hspace*{2em}令民为什伍。而相牧司连坐。不告奸者腰斩。告奸者与斩敌首同赏。匿奸者与降敌同罚。民有二男以上不分异者。倍其赋。有军功者。各以率受上爵。为私斗者。各以轻重被刑大小。僇力本业。耕织致粟帛多者复其身。事末利及怠而贫者。举以为收孥。宗室非有军功论。不得为属籍。明尊卑爵秩等级。各以差次名田宅。臣妾衣服以家次。有功者显荣。无功者虽富无所芬华。\\
\hspace*{2em}令既具。未布。恐民之不信。已乃立三丈之木於国都市南门。募民有能徙置北门者予十金。民怪之。莫敢徙。复曰﹁能徙者予五十金﹂。有一人徙之。辄予五十金。以明不欺。卒下令。\\
\hspace*{2em}令行於民期年。秦民之国都言初令之不便者以千数。於是太子犯法。卫鞅曰。﹁法之不行。自上犯之。﹂将法太子。太子。君嗣也。不可施刑。刑其傅公子虔。黥其师公孙贾。明日。秦人皆趋令。行之十年。秦民大说。道不拾遗。山无盗贼。家给人足。民勇於公战。怯於私斗。乡邑大治。秦民初言令不便者有来言令便者。卫鞅曰﹁此皆乱化之民也﹂。尽迁之於边城。其後民莫敢议令。\\
\hspace*{2em}於是以鞅为大良造。将兵围魏安邑。降之。居三年。作为筑冀阙宫庭於咸阳。秦自雍徙都之。而令民父子兄弟同室内息者为禁。而集小乡邑聚为县。置令。丞。凡三十一县。为田开阡陌封疆。而赋税平。平斗桶权衡丈尺。行之四年。公子虔复犯约。劓之。居五年。秦人富彊。天子致胙於孝公。诸侯毕贺。\\
\hspace*{2em}其明年。齐败魏兵於马陵。虏其太子申。杀将军庞涓。其明年。卫鞅说孝公曰。﹁秦之与魏。譬若人之有腹心疾。非魏并秦。秦即并魏。何者。魏居领阸厄之西。都安邑。与秦界河而独擅山东之利。利则西侵秦。病则东收地。今以君之贤圣。国赖以盛。而魏往年大破於齐。诸侯畔之。可因此时伐魏。魏不支秦。必东徙。东徙。秦据河山之固。东乡以制诸侯。此帝王之业也。﹂孝公以为然。使卫鞅将而伐魏。魏使公子卬将而击之。军既相距。卫鞅遗魏将公子卬书曰。﹁吾始与公子驩。今俱为两国将。不忍相攻。可与公子面相见。盟。乐饮而罢兵。以安秦魏。﹂魏公子卬以为然。会盟已。饮。而卫鞅伏甲士而袭虏魏公子卬。因攻其军。尽破之以归秦。魏惠王兵数破於齐秦。国内空。日以削。恐。乃使使割河西之地献於秦以和。而魏遂去安邑。徙都大梁。梁惠王曰。﹁寡人恨不用公叔痤之言也。﹂卫鞅既破魏还。秦封之於。商十五邑。号为商君。\\
\hspace*{2em}商君相秦十年。宗室贵戚多怨望者。赵良见商君。商君曰。﹁鞅之得见也。从孟兰皋。今鞅请得交。可乎。﹂赵良曰。﹁仆弗敢原也。孔丘有言曰。﹁推贤而戴者进。聚不肖而王者退。﹂仆不肖。故不敢受命。仆闻之曰。﹁非其位而居之曰贪位。非其名而有之曰贪名。﹂仆听君之义。则恐仆贪位贪名也。故不敢闻命。﹂商君曰。﹁子不说吾治秦与。﹂赵良曰。﹁反听之谓聪。内视之谓明。自胜之谓强。虞舜有言曰。﹁自卑也尚矣。﹂君不若道虞舜之道。无为问仆矣。﹂商君曰。﹁始秦戎翟之教。父子无别。同室而居。今我更制其教。而为其男女之别。大筑冀阙。营如鲁卫矣。子观我治秦也。孰与五羖大夫贤。﹂赵良曰。﹁千羊之皮。不如一狐之腋。千人之诺诺。不如一士之谔谔。武王谔谔以昌。殷纣墨墨以亡。君若不非武王乎。则仆请终日正言而无诛。可乎。﹂商君曰。﹁语有之矣。貌言华也。至言实也。苦言药也。甘言疾也。夫子果肯终日正言。鞅之药也。鞅将事子。子又何辞焉。﹂赵良曰。﹁夫五羖大夫。荆之鄙人也。闻秦缪公之贤而原望见。行而无资。自粥於秦客。被褐食牛。期年。缪公知之。举之牛口之下。而加之百姓之上。秦国莫敢望焉。相秦六七年。而东伐郑。三置晋国之君。一救荆国之祸。发教封内。而巴人致贡。施德诸侯。而八戎来服。由余闻之。款关请见。五羖大夫之相秦也。劳不坐乘。暑不张盖。行於国中。不从车乘。不操干戈。功名藏於府库。德行施於後世。五羖大夫死。秦国男女流涕。童子不歌谣。舂者不相杵。此五羖大夫之德也。今君之见秦王也。因嬖人景监以为主。非所以为名也。相秦不以百姓为事。而大筑冀阙。非所以为功也。刑黥太子之师傅。残伤民以骏刑。是积怨畜祸也。教之化民也深於命。民之效上也捷於令。今君又左建外易。非所以为教也。君又南面而称寡人。日绳秦之贵公子。诗曰。﹁相鼠有体。人而无礼。人而无礼。何不遄死。﹂以诗观之。非所以为寿也。公子虔杜门不出已八年矣。君又杀祝懽而黥公孙贾。诗曰。﹁得人者兴。失人者崩。﹂此数事者。非所以得人也。君之出也。後车十数。从车载甲。多力而骈胁者为骖乘。持矛而操闟戟者旁车而趋。此一物不具。君固不出。书曰。﹁恃德者昌。恃力者亡。﹂君之危若朝露。尚将欲延年益寿乎。则何不归十五都。灌园於鄙。劝秦王显岩穴之士。养老存孤。敬父兄。序有功。尊有德。可以少安。君尚将贪商於之富。宠秦国之教。畜百姓之怨。秦王一旦捐宾客而不立朝。秦国之所以收君者。岂其微哉。亡可翘足而待。﹂商君弗从。\\
\hspace*{2em}後五月而秦孝公卒。太子立。公子虔之徒告商君欲反。发吏捕商君。商君亡至关下。欲舍客舍。客人不知其是商君也。曰。﹁商君之法。舍人无验者坐之。﹂商君喟然叹曰。﹁嗟乎。为法之敝一至此哉。﹂去之魏。魏人怨其欺公子卬而破魏师。弗受。商君欲之他国。魏人曰。﹁商君。秦之贼。秦彊而贼入魏。弗归。不可。﹂遂内秦。商君既复入秦。走商邑。与其徒属发邑兵北出击郑。秦发兵攻商君。杀之於郑黾池。秦惠王车裂商君以徇。曰。﹁莫如商鞅反者。﹂遂灭商君之家。\\
\hspace*{2em}太史公曰。商君。其天资刻薄人也。迹其欲干孝公以帝王术。挟持浮说。非其质矣。且所因由嬖臣。及得用。刑公子虔。欺魏将卬。不师赵良之言。亦足发明商君之少恩矣。余尝读商君开塞耕战书。与其人行事相类。卒受恶名於秦。有以也夫。\\

% ===== 索隐述赞（司马贞《史记索隐》）· 统一样式 =====
% 题头「索隐述赞曰」五字并入反白框；正文每四句一列，定宽盒 \hfil 居列中；
% \fontsize{24}{41} 列距对齐全书界行（121+41k）；\clearpage 使述赞独占一页。
\clearpage
\shuzanhead{索隐述赞曰}\linebreak
{\shuzanbody\fontsize{24}{41}\selectfont
\hbox to \hsize{\hfil 卫鞅入秦。景监是因。王道不用。霸术见亲。\hfil}\linebreak
\hbox to \hsize{\hfil 政必改革。礼岂因循。既欺魏将。亦怨秦人。\hfil}\linebreak
\hbox to \hsize{\hfil 如何作法。逆旅不宾。\hfil}%
\par}
\end{document}
